\section{Why the gate could not see it}
\label{sec:mechanism}

\subsection{The diversity was in the wrong place}

The two paths differ in language, data structure and intermediate
representation. They do not differ in where their meaning comes from. The Python
path imports its classification constants and file loaders from the
reconciliation module. The SPARQL path queries a graph built by a script that
imports the same constants and loaders from the same module. We computed the
overlap from the source at the published commit: the names each path imports
from that module are identical, four vocabularies and two loaders
(Appendix~\ref{app:artifacts}, \texttt{shared\_ancestor\_audit.py}).

\begin{center}\small
\begin{tabular}{l}
\texttt{reconcile.py} defines \texttt{GONE}, \texttt{LEFT\_EARTH},
\texttt{ERROR}, \texttt{LOST}, \texttt{load\_gcat}, \texttt{load\_celestrak} \\
$\downarrow$ \hfill $\downarrow$ \\
\texttt{governance\_report.py} $\rightarrow$ Python counts \hfill
\texttt{build\_graph.py} $\rightarrow$ RDF graph $\rightarrow$ SPARQL counts \\
\end{tabular}
\end{center}

The gate diversified the last step of the computation and shared the first.
Whatever those sets assert about the source vocabulary is asserted identically on
both sides, so the agreement the gate reports is evidence that two pieces of
code, given the same premises, drew the same conclusion. In one sentence: the
gate verified that the Python and SPARQL paths implemented the same
misunderstanding identically.

\subsection{The defective definition}

The pre-correction disposition rule was one line applied to a set literal. An
object counted as gone in GCAT if any of its status codes was in

\begin{quote}\small\ttfamily
GONE = \{"R","D","L","LF","S","F","AF","AS","AR","AR IN","AL","AL IN","TX"\}
\end{quote}

\noindent and, through \texttt{bool(st \& GONE) != (n in ctdec)}, every other
object counted as still in orbit. The error lived in that default. Codes the set
did not name, among them \texttt{E} (exploded), \texttt{C} (collided),
\texttt{DK} (docked) and \texttt{ATT} (attached), were silently read as ``in
orbit''. A rule that made a claim only on explicit membership of \texttt{GONE}
or of the in-orbit set, and made no claim otherwise, returns 208 on the frozen
data, and all 208 are confirmed by the full-history check of
Section~\ref{sec:second}. The missing codes mattered far less than the decision
about what to do with codes nobody had classified.

A smaller defect in the same module shows the pathology from another angle. GCAT
marks uncertain dates with a trailing question mark, so a status can arrive as
\texttt{"R?"}. The original code compared raw strings and failed to match it,
which moved 15 objects in or out of the count. None of these corrections is
discoverable by comparing outputs. Each is a correction to a reading of the
source.

\subsection{Relation to multiversion programming}

Independently produced program versions do not fail independently
\citep{knight1986experimental}, and the accompanying theory located the cause in
the input rather than the versions: some inputs are hard for everyone, so
failures correlate however separately the versions were written
\citep{eckhardt1985theoretical, littlewood1989conceptual}. N-version designs were
proposed precisely to buy independence through separate development
\citep{avizienis1985nversion}. Our case differs in one structural respect that
makes it worse. In multiversion programming the shared artifact is a
specification, and versions can at least diverge in how they read it. Here the
shared artifact is an executable set of constants imported by both paths, so
there is no input on which they could disagree about a status code. The
correlation is not statistical but total.
